\documentclass[11pt,a4paper]{article}
\usepackage[utf8]{inputenc}
\usepackage[vietnamese]{babel}
\usepackage{amsmath,amssymb}
\usepackage{geometry}
\geometry{a4paper, total={180mm,260mm}, left=15mm, right=15mm, top=18mm, bottom=18mm}

\usepackage{tikz}
\usetikzlibrary{shapes.geometric, arrows.meta, positioning, calc}
\tikzset{
  arr/.style={-{Stealth[scale=0.8]}},
  every node/.append style={align=center}
}
\usepackage{tcolorbox}
\tcbuselibrary{skins,breakable}
\usepackage{fontawesome}
\usepackage{fancyhdr}
\usepackage{enumitem}

\pagestyle{fancy}
\fancyhf{}
\renewcommand{\headrulewidth}{0.5pt}
\renewcommand{\footrulewidth}{0.5pt}
\fancyhead[L]{\small\bfseries\color{violet!85!black}\faBook\ SỔ TAY KIẾN THỨC TOÁN 2}
\fancyhead[R]{\small\bfseries\color{violet!85!black}Lớp Toán Thầy Lĩnh -- loptoanthaylinh.vn}
\fancyfoot[C]{\small\color{gray}-- Sổ tay kiến thức: Bài 25. Điểm, đoạn thẳng, đường thẳng, đường cong --}

\setlength{\parindent}{0pt}
\setlength{\parskip}{0pt}

\begin{document}

% --- BANNER BÀI HỌC ---
\begin{tcolorbox}[
  enhanced,
  colback=violet!8!white,
  colframe=violet!70!black,
  arc=3mm, boxrule=1.2pt,
  left=3mm, right=3mm, top=2mm, bottom=2mm,
  halign=center
]
  {\LARGE\bfseries\color{violet!90!black}\faBookmark\ BÀI 25: ĐIỂM, ĐOẠN THẲNG, ĐƯỜNG THẲNG, ĐƯỜNG CONG}\\[1mm]
  {\itshape\footnotesize\color{black!75} Nhận biết các yếu tố hình học cơ bản: Điểm, đoạn thẳng, đường thẳng, đường cong và ba điểm thẳng hàng trên 1 trang A4}
\end{tcolorbox}

\vspace{-2mm}

\noindent
\begin{minipage}[t]{0.49\textwidth}\vspace{0pt}
\begin{tcolorbox}[
  enhanced,
  colback=violet!4!white,
  colframe=violet!75!black,
  arc=2mm, boxrule=0.8pt,
  left=2mm, right=2mm, top=1.8mm, bottom=1.8mm,
  title={\footnotesize\bfseries\color{white}\faKey\ 1. Điểm, đoạn thẳng, đường thẳng}
]
{\footnotesize
\begin{itemize}\setlength{\itemsep}{1pt}\setlength{\leftskip}{-2.5mm}
  \item \textbf{Điểm:} Một dấu chấm nhỏ trên mặt giấy, đặt tên bằng chữ cái in hoa: điểm $A$, điểm $B$.
  \item \textbf{Đoạn thẳng $AB$:} Nối hai điểm $A$ và $B$ bằng thước thẳng. Đoạn thẳng bị \textbf{giới hạn ở hai đầu} bởi hai điểm mút $A$ và $B$, có độ dài đo được cụ thể.
  \item \textbf{Đường thẳng $AB$:} Kéo dài đoạn thẳng $AB$ về \textbf{hai phía không bị giới hạn}.
  \item \textbf{Đường cong:} Đường uốn lượn, không thẳng.
  \item \textbf{Ba điểm thẳng hàng:} Khi ba điểm cùng nằm trên một đường thẳng.
\end{itemize}
}
\end{tcolorbox}
\end{minipage}
\hfill
\begin{minipage}[t]{0.49\textwidth}\vspace{0pt}
\begin{tcolorbox}[
  enhanced,
  colback=teal!4!white,
  colframe=teal!75!black,
  arc=2mm, boxrule=0.8pt,
  left=2mm, right=2mm, top=1.8mm, bottom=1.8mm,
  title={\footnotesize\bfseries\color{white}\faSliders\ 2. Kĩ năng nhận biết và kiểm tra}
]
{\footnotesize
\begin{itemize}\setlength{\itemsep}{1pt}\setlength{\leftskip}{-2.5mm}
  \item \textbf{Đo độ dài đoạn thẳng:}
  \begin{itemize}\setlength{\itemsep}{0.5pt}\setlength{\leftskip}{-2.5mm}
    \item Đặt vạch $0$ của thước trùng với điểm đầu.
    \item Điểm cuối trùng với vạch số nào thì đó là độ dài đoạn thẳng.
  \end{itemize}
  \item \textbf{Kiểm tra 3 điểm thẳng hàng:}
  Đặt mép thước thẳng đi qua hai điểm đầu, nếu điểm thứ ba cũng nằm sát mép thước thì ba điểm đó \textbf{thẳng hàng}.
  \item \textbf{Phân biệt:} Đoạn thẳng có 2 điểm chặn ở đầu mút; đường thẳng kéo dài mãi mãi.
\end{itemize}
}
\end{tcolorbox}
\end{minipage}

\vspace{-1.5mm}
\begin{tcolorbox}[
  enhanced,
  colback=orange!4!white,
  colframe=orange!80!black,
  arc=2mm, boxrule=0.8pt,
  left=2mm, right=2mm, top=1.5mm, bottom=1.5mm,
  title={\footnotesize\bfseries\color{white}\faSitemap\ 3. Sơ đồ hình vẽ trực quan 4 đối tượng hình học cốt lõi}
]
\centering
\begin{tikzpicture}[x=1.3cm, y=0.8cm]
  % Điểm
  \filldraw[violet!90!black] (0, 0.5) circle (2.5pt) node[below=2pt, font=\scriptsize\bfseries] {Điểm $A$};

  % Đoạn thẳng
  \draw[thick, teal!80!black] (1.5, 0.5) -- (3.5, 0.5);
  \filldraw[teal!90!black] (1.5, 0.5) circle (2pt) node[below=2pt, font=\tiny\bfseries] {$M$};
  \filldraw[teal!90!black] (3.5, 0.5) circle (2pt) node[below=2pt, font=\tiny\bfseries] {$N$};
  \node[font=\scriptsize\bfseries, teal!90!black] at (2.5, 0.9) {Đoạn thẳng $MN$};

  % Đường thẳng
  \draw[thick, orange!85!black] (4.5, 0.5) -- (7.5, 0.5);
  \filldraw[orange!90!black] (5.2, 0.5) circle (2pt) node[below=2pt, font=\tiny\bfseries] {$C$};
  \filldraw[orange!90!black] (6.8, 0.5) circle (2pt) node[below=2pt, font=\tiny\bfseries] {$D$};
  \node[font=\scriptsize\bfseries, orange!90!black] at (6.0, 0.9) {Đường thẳng $CD$};

  % Đường cong
  \draw[thick, purple!80!black] (8.2, 0.2) to[out=45, in=135] (9.2, 0.8) to[out=-45, in=-135] (10.2, 0.2);
  \node[font=\scriptsize\bfseries, purple!90!black] at (9.2, 0.9) {Đường cong};
\end{tikzpicture}
\end{tcolorbox}

\vspace{-1.5mm}
\noindent
\begin{minipage}[t]{0.49\textwidth}\vspace{0pt}
\begin{tcolorbox}[
  enhanced,
  colback=green!4!white,
  colframe=green!60!black,
  arc=2mm, boxrule=0.8pt,
  left=2mm, right=2mm, top=1.8mm, bottom=1.8mm,
  title={\footnotesize\bfseries\color{white}\faCheckCircle\ 4. Ví dụ thực hành chuẩn mực}
]
{\footnotesize
\textbf{1. Đọc tên các đoạn thẳng:}\\
Cho tam giác $ABC$. Kể tên các đoạn thẳng:\\
$\implies$ Có $3$ đoạn thẳng: $AB, BC, CA$.
\textbf{2. Kiểm tra thẳng hàng:}\\
Cho 3 điểm $A, B, C$ cùng nằm trên đường thẳng d. Ta nói: Ba điểm $A, B, C$ \textbf{thẳng hàng}.\\
Điểm $B$ nằm giữa hai điểm $A$ và $C$.
}
\end{tcolorbox}
\end{minipage}
\hfill
\begin{minipage}[t]{0.49\textwidth}\vspace{0pt}
\begin{tcolorbox}[
  enhanced,
  colback=red!4!white,
  colframe=red!75!black,
  arc=2mm, boxrule=0.8pt,
  left=2mm, right=2mm, top=1.8mm, bottom=1.8mm,
  title={\footnotesize\bfseries\color{white}\faExclamationTriangle\ 5. Mẹo tránh bẫy trắc nghiệm}
]
{\footnotesize
\begin{itemize}\setlength{\itemsep}{1pt}\setlength{\leftskip}{-2.5mm}
  \item \textbf{Phân biệt Đoạn thẳng và Đường thẳng:}\\
  - Đoạn thẳng $AB$: Bị giới hạn ở 2 đầu $A$ và $B$, đo được độ dài.\\
  - Đường thẳng $AB$: Kéo dài vô tận về hai phía.
  \item Đặt tên điểm phải viết bằng chữ cái \textbf{in hoa} ($A, B, C$), không dùng chữ thường ($a, b$).
\end{itemize}
}
\end{tcolorbox}
\end{minipage}


\end{document}
